% Replacement for the C3 paragraph of Section V-C (the paragraph that currently introduces
% Fig. 6 with the analytical 7.36 -> 1.66 ms and 0.147 -> 0.033 ms numbers).
% All numbers below come from results/per_run.csv (30 seeds, 24,000 events per run, H = 0).

To examine how the payload reduction translates into upload latency when several clients share a
constrained uplink, we drive a trace-driven discrete-event simulation with the measured payload pairs.
Each of the 24{,}000 events carries its raw and its distilled representation, and both are submitted
with identical client assignments, arrival times and link conditions, so that only the payload
representation differs. Clients generate uploads as a Poisson process at one upload per second each,
and all uploads share a single FIFO uplink whose service time is $8(D_i+H)/R$ with $H=0$ in the
baseline.

Fig.~\ref{fig:upload_time}(a) reports the 95th percentile of the upload completion time, which includes
queueing and transmission, as the client population grows from ten to one hundred at 1~Mbps. At ten
clients, the P95 completion time is 12.6\,ms for raw episodes and 2.8\,ms for distilled experiences;
with one hundred clients, the raw P95 rises to 53.6\,ms while the distilled P95 stays at 3.4\,ms. The
gap therefore widens from $4.44\times$, which is exactly the $4.44\times$ payload ratio, to
$15.9\times$. The additional benefit comes from queueing rather than from the payload itself: at one
hundred clients the raw uploads drive the uplink to 74\% utilisation, whereas distilled uploads use
17\%.

Fig.~\ref{fig:upload_time}(b) isolates the load effect at a fixed population of fifty clients and
reports the reduction in P95 completion time. Raising the offered load generated by raw uploads from
0.30 to 0.95 increases the reduction from 83.5\% to 98.6\%, far beyond the 77.5\% that the payload ratio
alone would give; over the same range the raw P95 grows from 18.1\,ms to 264.2\,ms while the distilled
representation stays within 3.0--3.6\,ms. At low utilisation the simulated mean completion time recovers
the analytical baseline of Eq.~(2) (7.71\,ms against 7.36\,ms for raw episodes and 1.67\,ms against
1.66\,ms for distilled experiences), so the divergence at higher load is attributable to queueing.
Adding 60\,bytes of per-upload protocol overhead raises the P95 completion time at fifty clients by
2.1\,ms for raw episodes and 0.7\,ms for distilled experiences without changing the ordering. These
results indicate that payload reduction shortens the transmission of each upload and also removes
queueing delay, and that its benefit grows as uplink contention increases.
